
%+PATCH,//BRILL/TEX
%+DECK,examples,T=LATEX.

\renewcommand*\thesection{\Roman{section}}
\section{Examples}

There exsists  in the \gti{input} directory a file for every given example.

You may copy this file to

\cbox{.pybrill\_last.dat}

in the working directory, start \pB and trigger the calculations
by pressing \gti{Brilliance/Calculate} or \gtie{Spectra/Calculate}.
However, it is recommended to start with the first example and use
\git{Set-Up} to preceed to the next example.

\subsection{Basic Beam Electron Monte-Carlo}

\begin{figure}[htb]
   \centering
   \includegraphics[width=0.55\textwidth]
   {pictures/ex1\_setup.png}
   \vspace*{-0.2cm}
   \caption{\gti{Set-Up/Spectra} for example 1}
   \label{ex1fig1}
\end{figure}

Before we start, we clean up the working directory to have a well defined, reproducable
setting:

\cbox{
rm .pybrill\_last.dat\\
rm urad\_phase.bun\\
rm urad\_phase.fld\\
rm urad\_phase.flx\\
rm urad\_phase.geo\\
rm urad\_phase.nor\\
rm urad\_phase.pin\\
rm urad\_phase.run
}

We start \pBe, press \gti{Set-Up}, and set:

\cbox{
Nelec = 1000\\
Number of Photon Energies = 11\\
Min. Photon Energy = 695.0\\
Max. Photon Energy = 705.0 \\
Width of Pinhole = 1.0 \\
Height of Pinhole = 1.0
}

We close the menu and click

\cbox{Spectra/Calculate}

We see, that the calculations take about a minute, depending on your
computer.

\cbox{Spectra/Plot/Flux}

\begin{figure}[htb]
   \centering
   \includegraphics[width=0.5\textwidth]
   {pictures/ex1\_flx.pdf}
   \vspace*{-0.2cm}
   \caption{\gti{Flux spectrum}}
   \label{ex1fig2}
\end{figure}

The flux spectrum we see, stems from 1000 electrons randomly chosen in
there phase-space. Thus, the result will probably change, when we
rerun the calculations. However, in the set-up you find an option to
fix the starting seed of the random number generator. This allows to
reproduce results of a program run.

We have a look on the flux-density distribution in the pinhole (Fig.
\ref{ex1fig2}:

\cbox{Spectra/Plot/Distribution in Pinhole/Stokes/S0}

\begin{figure}[htb]
   \centering
   \includegraphics[width=0.5\textwidth]
   {pictures/ex1\_scat2d.pdf}
   \vspace*{-0.2cm}
   \caption{\gti{Flux spectrum, plotted with option 'scat2d'}}
   \label{ex1fig3}
\end{figure}

\begin{figure}[htb]
   \centering
   \includegraphics[width=0.5\textwidth]
   {pictures/ex1\_inter.pdf}
   \vspace*{-0.2cm}
   \caption{\gti{Flux spectrum, plotted with option 'inter'}}
   \label{ex1fig4}
\end{figure}

We change the plotting options by selecting:

\cbox{Spectra/Plotting Options}

and choosing 'inter' and plot the distribution again
(Fig. \ref{ex1fig3}
).

\subsection{Combining Beam Electron and Observation Point Monte-Carlo}

For a large number of electrons and many grid points, the calculation
time can increase drastically. To speed up the calculations of the
flux through the pinhole,
we select in \gti{Set-Up/Spectra}

\cbox{
Nelec = 100000\\
Each n\_th electron is recorded = 100 \\
Number of Photons Energies = 1}

We close the menu and start the calculations.

\begin{figure}[htb]
   \centering
   \includegraphics[width=0.5\textwidth]
   {pictures/ex2\_flux.pdf}
   \vspace*{-0.2cm}
   \caption{\gti{Flux through pinhole with one sigma confidence level}}
   \label{ex2flux}
\end{figure}

Now, we only have to wait some seconds due to

\cbox{Number of Photons Energies = 1}

In this mode, for each electron a single observation point is
randomly choosen in the pinhole. We see that flux spectrum shows only small statistical
uncertainty (Fig. \ref{ex2flux}).

\begin{figure}[htb]
   \centering
   \includegraphics[width=0.5\textwidth]
   {pictures/ex2\_fdpin.pdf}
   \vspace*{-0.2cm}
   \caption{\gti{Flux-density in pinhole, plotted with option 'boxes'}}
   \label{ex2fdpin}
\end{figure}

However, when we look at the distribution of the flux-density in the
pinhole, we see clearly fluctions due to this special Monte-Carlo mode (Fig.
\ref{ex2fdpin}).

But in many cases, only the spetra of the flux, the flux-density, or
the profiles of the distributions (Fig. \ref{ex2fdprof}) are of interest.

\begin{figure}[htb]
   \centering
   \includegraphics[width=0.5\textwidth]
   {pictures/ex2\_fdprof.pdf}
   \vspace*{-0.2cm}
   \caption{\gti{Profile of flux-density in pinhole}}
   \label{ex2fdprof}
\end{figure}

Also the reduced number of sampled
electrons speeds up the calculations, since the I/O is reduced.

\begin{figure}[htb]
%   \vspace*{-2.5cm}
   \centering
   \includegraphics[width=0.5\textwidth]
   {pictures/ex2\_zzp.pdf}
   \vspace*{-0.2cm}
   \caption{
   \gti{Horizontal phase-space of beam electrons sample}}
   \label{ex2zzp}
\end{figure}

To check the phase-space a sample of 1000 electrons is certainly
sufficient (Fig. \ref{ex2zzp}).

\subsection{Back Propagation of the Radiation to the Source Plane}

To calculate the source size, we propagate the radiation back to the
source plane by selecting in \gti{Set-Up/Spectra}\\

\cbox{
Nelec = 1\\
Number of Photons Energies = 0 \\
Width of Pinhole = 2.0 \\
Height of Pinhole = 2.0 \\
Number of hori. points = 101 \\
Number of vert. points = 101 \\
Propagate radiation field back to origin = 1 \\
Width of Pinhole in Plane = 0.1 \\
Height of Pinhole in Plane = 0.1 \\
Number of hori. points in Plane = 51 \\
Number of vert. points in Plane = 51
}\\

and starting the calculations. Then we switch on the statistic option in
the menu

\cbox{Spectra/Plotting options}

and plot the horizontal cut of the Stokes
$S_0$ in the menu

\cbox{Plot/Distributions of Propagated Fields}

\begin{figure}[htb]
   \centering
   \includegraphics[width=0.5\textwidth]
   {pictures/ex3\_pinhcut.pdf}
   \vspace*{-0.2cm}
   \caption{
   \gti{Horizontal cut of $S_0$ in observation plane}}
   \label{ex3pinh}
\end{figure}

For the product of the RMS values of the two plots

\begin{eqnarray*}
\sigma_r \times \sigma_{r'} &=&
\frac{\lambda}{4 \pi}
\end{eqnarray*}
	
we find

\begin{figure}[htb]
   \centering
   \includegraphics[width=0.5\textwidth]
   {pictures/ex3\_prophcut.pdf}
   \vspace*{-0.2cm}
   \caption{
   \gti{Horizontal cut of $S_0$ in source plane}}
   \label{ex3proph}
\end{figure}

\begin{eqnarray*}
0.1074 ~mm \times 0.01846 ~mm / 10~m &=&
0.1982604 ~nm
\end{eqnarray*}

We have $\lambda = 1.77 ~nm$ for $700 ~eV$, thus

\begin{eqnarray*}
\frac{1.77~nm}{4 \pi} &=& 0.141 ~ nm
\end{eqnarray*}

which underestimates the numerical value by a factor of $\sqrt{2}$.

\subsection{Generating Photons from Undulator Source for Ray-tracing}

\PyB uses the cross product of the calculated electric and magnetic
fields to generate photons as input for ray-tracing codes.

\begin{figure}[htb]
   \centering
   \includegraphics[width=0.5\textwidth]
   {pictures/ex4\_ztz.pdf}
   \vspace*{-0.2cm}
   \caption{
   \gti{Horizontal phase-space of generated photons}}
   \label{ex4ztz.pdf}
\end{figure}

Due to the applied algorithm
in this examples the photons are generated for an exact beam energy,
i.e. without energy-spread.

To generate photons for electrons with energy-spread the program must
sample the beam energy in discrete steps. The number of steps can
be set in the set-up menu.

Note, if you run the program for another beam energy by hand, to see
the effect of beam energy deviations, make shure that the undulater is
set-up such that the fields are not adjusted according to the actual
energy,
as it is done, when a fixed harmonic is selected in the undulator
set-up menu.

